% Run GCT_KRON_20261120; GPT-6 Astra Pro; effort not exposed.
\section{Grouped derivatives give a cleaner universal degree gate}

\subsection*{A-023. Improving the recurrence by treating each derivative power at once}
For $L\geq1$ and $a\geq2$, define
\[
 P_2(L)=1,\qquad s=\min(a-1,L),\qquad
 P_a(L)=s+\max\{P_{a-1}(L),\min((a-2)s,L)\}
                        \quad(a\geq3).                \tag{A-023.1}
\]
Then for every $\lambda\vdash ab$ with $\lambda_2\leq L$,
\[
 b\geq P_a(L)
 \quad\Longrightarrow\quad
 I(\operatorname{Ch}_b(V))_a[\lambda]=0                 \tag{A-023.2}
\]
over every finite-dimensional complex $V$. Moreover
\[
           P_a(L)\leq N_a(L),\qquad P_a(L)\leq(a-1)^2. \tag{A-023.3}
\]
In particular $b\geq(a-1)^2$ gives the uniform vanishing
in every type. These are sufficient bounds, not claims of
optimal stabilization.

For the padded permanent comparison with $n\geq m\geq1$
and outer degree $d\geq2$, there is no actual degree-$d$
multiplicity obstruction whenever
\[
                              n\geq P_d(dm).           \tag{A-023.4}
\]
For $L\geq d$, the exact threshold is
\[
 \boxed{\quad P_d(L)=\binom d2+
       \max\left(0,\max_{2\leq j\leq d-1}
           \min\left\{\binom j2,L-\binom j2\right\}\right).\quad}
                                                               \tag{A-023.5}
\]
At $d=2$ the maximum is zero. Since the inner minimum is at
most $\lfloor L/2\rfloor$, a simple sufficient condition is
\[
       \boxed{\quad n\geq
              \left\lfloor\frac{d(d+m-1)}2\right\rfloor.\quad}
                                                               \tag{A-023.6}
\]
The exact maximum can improve this further. Thus an actual
multiplicity obstruction necessarily satisfies
\[
 n<P_d(dm)\leq
 \min\left\{(d-1)^2,
             \left\lfloor\frac{d(d+m-1)}2\right\rfloor\right\}.
                                                               \tag{A-023.7}
\]

\paragraph{The power of a differential operator.}
Let $D=D_{Y,U}$ on polynomial functions homogeneous of degrees
$p,q$ in independent vector blocks $U,Y$, and let $e$ be an
integer with $1\leq e\leq p$. With arbitrary polynomial
spectators of fixed degrees, $D^e$ is injective on type
$\lambda$ whenever
\[
                         p-e\geq\min(q,\lambda_2).     \tag{A-023.8}
\]
This is stronger than requiring every individual derivative to
be injective on its entire intermediate homogeneous space.

\begin{proof}
Use A-020's multiplicity-free Pieri decomposition of the two
blocks. On the summand $S_{(p+q-i,i)}V^*$, its highest-weight
vector is
$h_i=(U_1Y_2-U_2Y_1)^iU_1^{p-i}Y_1^{q-i}$.
Applying $D^e$ multiplies by the falling factorial
$(p-i)(p-i-1)\cdots(p-i-e+1)$ and changes the exponents
of $U_1,Y_1$ by $-e,+e$. It is nonzero exactly for
$i\leq p-e$. Thus the kernel is the sum of the summands
with $p-e<i\leq\min(p,q)$.

If $q\leq p-e$, this kernel is zero. Otherwise every
constituent of the kernel has second row at least $p-e+1$;
tensoring the spectator symmetric powers can only add boxes
by horizontal Pieri. It therefore has no type $\lambda$
if $\lambda_2\leq p-e$. These two cases prove (A-023.8).
The one-dimensional case follows from the same monomial
derivative, and types not fitting the dimension are zero.
\end{proof}

\begin{proof}[Proof of the refined recurrence]
The base $a=2$ and one-row types are the same as in A-021.
Fix an upper bound $L$ on the second row, let $r=a-1$,
$s=\min(r,L)$, and $m=b-s$. By (A-023.1),
\[
                m\geq P_r(L)\geq1,
                \qquad m\geq\min((r-1)s,L).            \tag{A-023.9}
\]
Let $T$ be in the $\lambda$ isotypic component of symmetric
$a$-linear forms polarizing the degree-$a$ Chow ideal.
The binary-factorization identity and the first differential
sequence of A-021 apply without change: at each step the
source degree is at least $s+1$, while the untouched
destination degree is $r$; either $s=r$ gives $p>q$, or
$s=L$ gives $p>L\geq\lambda_2$. Hence
$T(x^sQ,\ldots,x^sQ,t^b)=0$.

Retaining the $x,t$ polynomial variables gives the image in
$\operatorname{Sym}^{rs}V^*\otimes\operatorname{Sym}^bV^*
\otimes I(\operatorname{Ch}_m(V))_r$.
The induction hypothesis kills all its last-factor types
whose second row is at most $L$. The two symmetric-power
spectators cannot lower the second row, so the entire
$\lambda$ isotypic image is zero. Consequently
\[
 T(x^sQ_1,\ldots,x^sQ_r,t^b)=0
                   \quad(Q_i\in\operatorname{Sym}^mV). \tag{A-023.10}
\]

Now start from $f=T(v_1^b,\ldots,v_r^b,t^b)$, and apply
the complete powers $D_{x,v_i}^{\,s}$ for $i=1,\ldots,r$.
Before the $i$-th power, its source block has degree $p=b$
and the common destination $x$ has degree $q=(i-1)s$.
The earlier powers affected only other source blocks.
By (A-023.9),
\[
 p-s=m\geq\min((r-1)s,L)\geq\min(q,\lambda_2).
\]
The grouped injectivity criterion (A-023.8) therefore applies
to each complete power on type $\lambda$. Their composition
is injective, while its value is
\[
 \left(\frac{b!}{m!}\right)^r
                   T(x^sv_1^m,\ldots,x^sv_r^m,t^b)=0
\]
by (A-023.10). Thus $f=0$, and pure-power spanning yields
$T=0$. This proves (A-023.2), retaining all multiplicity copies.

The new second term $\min((r-1)s,L)$ is at most
$\min(rs-1,L)$, since $s\geq1$. Induction therefore
gives $P_a(L)\leq N_a(L)$. For the uniform bound, use
$s\leq r$, $P_r(L)\leq(r-1)^2$, and
$\min((r-1)s,L)\leq r(r-1)$ in (A-023.1). This gives
$P_a(L)\leq r+r(r-1)=r^2=(a-1)^2$.
The base $P_2=1$ verifies both inductions.
\end{proof}

\begin{proof}[Proof of the explicit degree gate]
The padding morphism and Pieri argument in A-022 show that
every type occurring on the padded permanent has
$\lambda_2\leq dm$. Apply (A-023.2) with $a=d$, $b=n$,
and $L=dm$. If $n\geq P_d(dm)$ its Chow-ideal component
vanishes. Chow products are diagonal determinants, so the
determinant-ideal component also vanishes, and its actual
multiplicity equals the ambient multiplicity. This proves
(A-023.4) and the resulting comparison with every
permanent-side copy.

If $L\geq d$, every $s$ in the recursion is $a-1$.
The same unrolling as in A-022 gives
\[
 P_d(L)=\binom d2+
 \max\left(0,\max_{2\leq j\leq d-1}
       \left\{\min(j(j-1),L)-\binom j2\right\}\right).
\]
Writing $t=\binom j2$ turns the expression in braces
into $\min(2t,L)-t=\min(t,L-t)$, proving (A-023.5).
Its integer value is at most $\lfloor L/2\rfloor$;
the outer maximum with zero obeys the same bound for $L\geq1$.
Thus
\[
 P_d(dm)\leq\binom d2+\left\lfloor\frac{dm}{2}\right\rfloor
              =\left\lfloor\frac{d(d+m-1)}2\right\rfloor.
\]
This proves (A-023.6)--(A-023.7).

For exact evaluation at $d\geq3$, let
\[
 k=\left\lfloor\frac{1+\sqrt{1+4L}}2\right\rfloor,
\]
the largest integer satisfying $k(k-1)\leq L$.
For $j\leq k$ the summand $\min(\binom j2,L-\binom j2)$
is $\binom j2$ and increases. For $j\geq k+1$ it is
$L-\binom j2$ and decreases. Hence it is enough to inspect
$\min(d-1,k)$ and $\min(d-1,k+1)$. As $L\geq d\geq3$,
both clipped indices are at least two. The integer $k$ can
equivalently be computed as
$(1+\operatorname{isqrt}(1+4L))\mathbin{//}2$, where
\(\operatorname{isqrt}\) and \(//\) denote integer square
root and integer division. This gives an exact computation,
with no numerical root approximation. \end{proof}

\paragraph{Why grouping changes the bound.}
The power $D^s$ need only be injective on its actual source.
Some intermediate homogeneous spaces contain extra irreducible
summands that are never reached from that source. Requiring
injectivity on all of those larger spaces introduced the
extra $s-1$ in A-022's final derivative test. The explicit
Pieri calculation removes that loss. The proof still uses
additional structure beyond an effective $K_0$ class: the
equivariant evaluation and differential maps, their kernels,
and preservation of all multiplicity copies.

These are morphisms of polynomial comodules in the language
of SGA 3, I, Propositions 4.7.2--4.7.3, pp.~34--36, and
symmetric/Schur operations in SGA 6, VI, Theorem 3.3,
pp.~392--393. The kernel calculation uses horizontal Pieri,
Macdonald, I, (5.16), p.~73. As in A-020--A-022, the
underlying maps are integral while the irreducible and
nonzero-coefficient argument is made in characteristic zero.
The improvement concerns the actual multiplicity comparison;
it asserts no Kronecker saturation or conclusion about
$\mathsf P$ versus $\mathsf{NP}$.
